\chapter{Linear Algebra}
A unimodular matrix has determinant $\pm 1$


$$\rm{det}(\bf{AB}) = \rm{det}(\bf{A})\rm{det}(\bf{B})$$
The trace is invariant under cyclic permutations
\begin{align}
	\textrm{tr}(ABC) = \textrm{tr}(CAB) = \textrm{tr}(BCA)
\end{align}

$$a_i b_{ij} c_j = \textbf{a}\cdot\textbf{B}\cdot\textbf{b} = \sum_{i} a_i \sum_{j} b_{ij} c_j$$

\section{Tensors}
The rank of a tensor is simply the number of indices needed to describe it. A vector is a tensor of the first rank, and a general $n \times m$ matrix is a tensor of the second rank. The tensor product of two 2x2 matrices is given by
\begin{align}
\begin{bmatrix}
a_{11} & a_{12} \\
a_{21} & a_{22}
\end{bmatrix}
\otimes
\begin{bmatrix}
b_{11} & b_{12} \\
b_{21} & b_{22}
\end{bmatrix} &=
\begin{bmatrix}
a_{11} \begin{bmatrix}
b_{11} & b_{12} \\
b_{21} & b_{22}
\end{bmatrix}& a_{12} \begin{bmatrix}
b_{11} & b_{12} \\
b_{21} & b_{22}
\end{bmatrix}\\
a_{21} \begin{bmatrix}
b_{11} & b_{12} \\
b_{21} & b_{22}
\end{bmatrix}& a_{22}\begin{bmatrix}
b_{11} & b_{12} \\
b_{21} & b_{22}
\end{bmatrix}
\end{bmatrix}
\end{align}
We can do this for any vector space by generalizing the form above. For instance

\begin{align}
\begin{bmatrix}
a_{1}  \\
a_{2} 
\end{bmatrix}
\otimes
\begin{bmatrix}
b_{1}\\
b_{2}
\end{bmatrix} = 
\begin{bmatrix}
a_{1}  \begin{bmatrix}
b_{1}\\
b_{2}
\end{bmatrix} \\
a_{2} \begin{bmatrix}
b_{1}\\
b_{2}
\end{bmatrix} 
\end{bmatrix}
\end{align}

Another important tensor operator is the direct sum, which is written as
\begin{align}
	\begin{bmatrix}
a_{11} & a_{12} \\
a_{21} & a_{22}
\end{bmatrix}
\oplus
\begin{bmatrix}
b_{11} & b_{12} \\
b_{21} & b_{22}
\end{bmatrix} &=
\begin{bmatrix}
a_{11} & a_{12} & 0 & 0\\
a_{21} & a_{22} & 0 & 0\\
0 & 0& b_{11} & b_{12}\\
0 & 0 & b_{21} & b_{22}
\end{bmatrix}
\end{align}

\section{Determinant}
For an arbitrarily sized matrix, can iterate this procedure.
\begin{enumerate}
\item Write out each value in the top row, alternating sign ($+,-,+,-,...$)
\item Multiply each respective value by the determinant of the square matrix created by the rows and columns that the value does not belong to
\end{enumerate}
As an example a 4 x 4
\begin{align}
\begin{vmatrix} a & b & c & d\\e & f & g & h\\i & j & k & l\\m & n & o & p \end{vmatrix}=a\,\begin{vmatrix} f & g & h\\j & k & l\\n & o & p \end{vmatrix}-b\,\begin{vmatrix} e & g & h\\i & k & l\\m & o & p \end{vmatrix}+c\,\begin{vmatrix} e & f & h\\i & j & l\\m & n & p \end{vmatrix}-d\,\begin{vmatrix} e & f & g\\i & j & k\\m & n & o \end{vmatrix}
\end{align}
The matrix of a 3 x 3 is 
\begin{align}
 \begin{vmatrix} a & b & c\\d & e & f\\g & h & i \end{vmatrix} &= a\,\begin{vmatrix} e & f\\h & i \end{vmatrix} - b\,\begin{vmatrix} d & f\\g & i \end{vmatrix} + c\,\begin{vmatrix} d & e\\g & h \end{vmatrix}\end{align}
 And 2 x 2
 \begin{align}\begin{vmatrix} a & b\\c & d \end{vmatrix}=ad - bc .\end{align}

\subsection{Properties of the Determinant}
The determinant has several useful properties:
\begin{itemize}
    \item $\det(\textbf{A}\textbf{B}) = \det(\textbf{A})\det(\textbf{B})$
    \item $\det(\textbf{A}^T) = \det(\textbf{A})$
    \item $\det(\lambda \textbf{A}) = \lambda^n \det(\textbf{A})$ for an $n \times n$ matrix
    \item A matrix is invertible iff its determinant is non-zero
\end{itemize}

\section{Matrix Diagonalization}\label{diagonalize}
One can follow the procedure outlined below to diagonalize a matrix\cite{griffiths_qm}

\begin{enumerate}
\item Find eigenvalues ($\lambda_1, \lambda_2, ... , \lambda_N$) of matrix $\bf{M}$ with det($\bf{M}-\lambda\bf{I}$) = 0
\item Using the eigenvalues, find the corresponding eigenvectors by plugging in each to $(\bf{M}-\lambda_i\bf{I})\bf{a_i} = 0$
\item The diagonalized matrix $\bf{T}$ is now 

$$\bf{T} = \left(
{\begin{array}{cccc}
\lambda_1 & 0 &...&0 \\
0 & \lambda_2 & ...&0\\
\vdots & \vdots &\ddots & \vdots \\
0 & 0 & ... &\lambda_N
\end{array}}
\right)
$$

\item We need to find$$ \bf{M} = \bf{STS^{-1}}$$
\item $\bf{S^{-1}}$ is equal to each eigenvector laid out as a column of the matrix, placed according to where you place it with $\bf{T}$

$$\bf{S^{-1}} = \left(
{\begin{array}{cccc}
a(\lambda_1)_1 & a(\lambda_2)_1 &...&a(\lambda_N)_1 \\
a(\lambda_1)_2 & a(\lambda_2)_2 & ...&a(\lambda_N)_2\\
\vdots & \vdots &\ddots & \vdots \\
a(\lambda_1)_N & a(\lambda_2)_N & ... &a(\lambda_N)_N
\end{array}}
\right)
$$

\item $\bf{S}$ can be calculated by simply inverting the matrix $\bf{S^{-1}}$
\end{enumerate}

\section{Simultaneous Eigenstates}
To find a simultaneous set of eigenvectors for two operators, follow this:
\begin{itemize}
\item Given two 3x3 matrices

$$\bf{M} = 
\left(
{\begin{array}{ccc}
1&0&1\\
0&0&0\\
1&0&1\\
\end{array}}
\right)
~~~~\bf{D} = 
\left(
{\begin{array}{ccc}
2&1&1\\
1&0&-1\\
1&-1&2\\
\end{array}}
\right)$$

\item First solve for their eigenvalues ($\lambda_M = 0,0,2$, $\lambda_D = -1,2,3$)
\item Next find both sets of unnormalized eigenvectors
$$
\lambda_M\rightarrow  
\left(
{\begin{array}{c}
a\\
b\\
-a
\end{array}}
\right),
\left(
{\begin{array}{c}
a\\
b\\
-a
\end{array}}
\right),
\left(
{\begin{array}{c}
a\\
0\\
-a
\end{array}}
\right)~~~~
\lambda_D\rightarrow
\left(
{\begin{array}{c}
0\\
0\\
0
\end{array}}
\right),\left(
{\begin{array}{c}
a\\
a\\
-a
\end{array}}
\right),
\left(
{\begin{array}{c}
a\\
0\\
-a
\end{array}}
\right)$$
\item Because $\textbf{D}$ is not degenerate, correspond each vector to another vector for degenerate $\textbf{M}$. 
\item Then just normalize the vectors as usual, and you have a set of eigenvectors for both matrices
\end{itemize}

\section{Useful Properties}
\begin{itemize}
    \item If two matrices commute and are diagonalizable, they can be simultaneously diagonalized.
    \item The trace of a matrix is the sum of its eigenvalues (counting multiplicity).
    \item A real symmetric matrix always has real eigenvalues and an orthonormal eigenbasis (Spectral Theorem).
\end{itemize}
